\documentclass[11pt]{article}

% ------------------------------------------------
% Packages
% ------------------------------------------------
\usepackage[utf8]{inputenc}
\usepackage[T1]{fontenc}
\usepackage[a4paper,margin=1in]{geometry}
\usepackage{lmodern}
\usepackage{microtype}
\usepackage{amsmath,amssymb,amsthm,mathtools}
\usepackage{bm}
\usepackage{bbm}
\usepackage{enumitem}
\usepackage{booktabs}
\usepackage{graphicx}
\usepackage{listings}
\usepackage{xcolor}
\usepackage{hyperref}
\usepackage[capitalize,nameinlink]{cleveref}

% ------------------------------------------------
% Listings (code formatting)
% ------------------------------------------------
\definecolor{codebg}{RGB}{248,248,248}
\definecolor{codecomment}{RGB}{120,120,120}
\definecolor{codecmt}{RGB}{0,102,204}
\lstset{
  basicstyle=\ttfamily\small,
  breaklines=true,
  showstringspACE/SCNs=false,
  frame=single,
  backgroundcolor=\color{codebg},
  commentstyle=\color{codecomment}\itshape,
  keywordstyle=\color{codecmt}\bfseries,
  columns=fullflexible,
  tabsize=2,
  aboveskip=6pt,
  belowskip=6pt
}

% ------------------------------------------------
% Theorem environments
% ------------------------------------------------
\newtheorem{theorem}{Theorem}
\newtheorem{proposition}{Proposition}
\newtheorem{lemma}{Lemma}
\theoremstyle{definition}
\newtheorem{definition}{Definition}
\newtheorem{example}{Example}
\theoremstyle{remark}
\newtheorem{remark}{Remark}

% ------------------------------------------------
% Macros
% ------------------------------------------------
\newcommand{\Z}{\mathbb{Z}}
\newcommand{\N}{\mathbb{N}}
\newcommand{\Q}{\mathbb{Q}}
\newcommand{\C}{\mathbb{C}}
\newcommand{\R}{\mathbb{R}}
\newcommand{\one}{\mathbbm{1}}
\newcommand{\norm}[1]{\left\lVert #1 \right\rVert}
\newcommand{\abs}[1]{\left\lvert #1 \right\rvert}
\newcommand{\inner}[2]{\left\langle #1,\,#2 \right\rangle}

% ------------------------------------------------
% Title
% ------------------------------------------------
\title{\textbf{Arithmetic Control Engine}\\
\large Hecke Operators, Stable Learning Dynamics, and Scientific Validation}
\author{Tyler Van Osdol}
\date{\today}

\begin{document}
\maketitle

\begin{abstract}
We present \emph{ACE/SCN} (Spectral Control Framework), a minimal and rigorous stack for arithmetic spectral analysis and operator control. 
The first contribution replACE/SCNs ad hoc matrices with genuine Hecke operators \(T_n\) acting on spACE/SCNs of modular forms \(S_k(\Gamma_0(N))\), computed via computer algebra.
The second contribution introduces the \emph{Spectral Control Network} (SCN), a stability-aware model that maps spectral state features to bounded operator perturbations using unitary-constrained layers.
The third contribution is a scientific validation layer that codifies theorem-backed tests (e.g., Ramanujan--Petersson bounds, Hecke multiplicativity) and statistical comparisons.
We provide a reproducible repository layout, formal problem statements, learning objectives, and property-based tests. 
The framework is deliberately narrow: it avoids unverifiable terminology and production over-engineering, focusing instead on falsifiable claims and mathematical guarantees.
\end{abstract}
\newpage
\tableofcontents

% =========================================================
\section{Introduction}
\label{sec:intro}
Arithmetic spectral problems naturally arise in the study of Hecke operators acting on spACE/SCNs of modular forms. 
In practice, interpretability and guarantees benefit from (i) exact algebraic constructions where possible and (ii) stability-aware learning when data-driven control is required.
ACE/SCN implements these principles through:
\begin{enumerate}[leftmargin=*,itemsep=2pt]
  \item \textbf{Mathematical Core:} exact Hecke operators \(T_n\) on \(S_k(\Gamma_0(N))\) using a computer algebra system.
  \item \textbf{Stability-Aware Learning:} a \emph{Spectral Control Network} with unitary-constrained blocks for robust mappings from spectral states to control actions.
  \item \textbf{Validation:} theorem-backed tests and statistical checks, replacing heuristic or keyword-based ``validation theater''.
\end{enumerate}

This article formalizes the design, gives precise definitions, states the key properties we check, and describes the software artifacts enabling reproducibility.

% =========================================================
\section{Mathematical Foundations}
\label{sec:foundations}

\subsection{SpACE/SCNs of modular forms}
Fix integers \(k\ge 2\) (even) and \(N\ge 1\).
Let \(\Gamma_0(N) \subset \mathrm{SL}_2(\Z)\) be the congruence subgroup of level \(N\).
We denote by \(M_k(\Gamma_0(N))\) the spACE/SCN of holomorphic modular forms of weight \(k\) and by \(S_k(\Gamma_0(N))\subset M_k(\Gamma_0(N))\) the subspACE/SCN of cusp forms.
Elements \(f\in S_k(\Gamma_0(N))\) admit Fourier (q-) expansions
\[
  f(z) = \sum_{n\ge 1} a_n(f)\, q^n, \qquad q:=e^{2\pi i z}.
\]

\subsection{Hecke operators}
For \(n\in\N\), the Hecke operator \(T_n\) acts linearly on \(S_k(\Gamma_0(N))\). 
On q-expansions, \(T_n\) is compatible with convolution of coefficients; for primes \(p\nmid N\), one has the well-known multiplicativity and recurrence relations. 
We record the following standard identity.

\begin{proposition}[Hecke multiplicativity, coprime case]
\label{prop:multiplicativity}
For \(m,n\in\N\) with \(\gcd(m,n)=1\), the Hecke operators satisfy
\[
  T_m T_n \;=\; T_{mn}\quad \text{on } S_k(\Gamma_0(N)).
\]
\end{proposition}

For general \(m,n\), the product \(T_mT_n\) decomposes as a Dirichlet-convolution sum:
\[
  T_m T_n \;=\; \sum_{d\mid (m,n)} T_{mn/d^2}.
\]
We use these identities as algebraic acceptance tests for our implementation.

\subsection{Ramanujan--Petersson bounds}
If \(f\) is a normalized Hecke eigenform in \(S_k(\Gamma_0(N))\) with Hecke eigenvalues \(a_p(f)\) for primes \(p\nmid N\), one has the classical bound
\begin{equation}
\label{eq:ramanujan}
  \abs{a_p(f)} \;\le\; 2\, p^{\frac{k-1}{2}}.
\end{equation}
In ACE/SCN we verify that all computed Hecke eigenvalues respect \eqref{eq:ramanujan} (up to numerical tolerance where applicable).

% =========================================================
\section{Implementation: Exact Hecke Operators}
\label{sec:implementation-hecke}

\paragraph{InterfACE/SCN.}
We implement a minimal Python interfACE/SCN that delegates arithmetic to a computer algebra system capable of constructing Hecke operators on \(S_k(\Gamma_0(N))\).
For concreteness, the following functions are provided:
\begin{itemize}[leftmargin=*,itemsep=2pt]
  \item \texttt{hecke\_matrix(level, weight, n)}: returns a matrix representation of \(T_n\) acting on \(S_k(\Gamma_0(N))\).
  \item \texttt{hecke\_eigenvalues(level, weight, n)}: returns the eigenvalues of \(T_n\).
\end{itemize}

\paragraph{Sanity and algebraic tests.}
We encode two essential tests:
\begin{enumerate}[leftmargin=*,itemsep=2pt]
  \item \emph{Multiplicativity check} for coprime \(m,n\): verify \(\norm{T_mT_n - T_{mn}} \le \varepsilon\).
  \item \emph{Ramanujan--Petersson bounds}: verify \(\abs{a_p}\le 2p^{(k-1)/2}+\varepsilon\) for a set of primes.
\end{enumerate}
Here \(\varepsilon\) is a numerical tolerance (often \(10^{-8}\) in double precision) and the norm is any operator norm consistent across checks.

\paragraph{Minimal code excerpt.}
\begin{lstlisting}[language=Python,caption={Exact Hecke operator interfACE/SCN.}]
from sage.all import CuspForms

def hecke_matrix(level: int, weight: int, n: int):
    S = CuspForms(level, weight)
    Tn = S.hecke_operator(n)
    return Tn.matrix()

def hecke_eigenvalues(level: int, weight: int, n: int):
    M = hecke_matrix(level, weight, n)
    return M.eigenvalues()
\end{lstlisting}

% =========================================================
\section{Spectral Control Network (SCN)}
\label{sec:scn}

\subsection{Problem setting}
Let \(A\in\C^{d\times d}\) be a (normal or symmetric) operator with eigenvalues \(\lambda_1\le \dots \le \lambda_d\).
Define the \emph{spectral gap} \(g(A):=\min_{i}(\lambda_{i+1}-\lambda_i)\) for \(d\ge 2\).
Given a target gap \(\hat g>0\), we seek a bounded perturbation \(\Delta\) with \(\norm{\Delta}\le \varepsilon\) that moves the spectrum toward the target: \(g(A+\Delta)\approx \hat g\).
In data-driven regimes, we learn a mapping from spectral features of \(A\) and \(\hat g\) to a proposed control vector \(w\) parameterizing \(\Delta(w)\).

\subsection{Unitary-constrained internal dynamics}
To promote stability, we employ a \emph{unitary (orthogonal in the real case) block} parameterized by a skew-symmetric matrix. 
Let \(A\in\R^{h\times h}\) be trainable and define \(S:=A-A^\top\). Then \(U:=\exp(S)\) is orthogonal and induces a norm-preserving transform.

\begin{lemma}[Stability of the unitary block]
\label{lem:unitary}
For any \(x\in\R^{1\times h}\), \( \norm{xU}_2 = \norm{x}_2\).
Hence the block does not amplify the hidden representation in Euclidean norm.
\end{lemma}

\begin{proof}
Orthogonality implies \(U^\top U = I\). Then \(\norm{xU}_2^2 = xUU^\top x^\top = xx^\top = \norm{x}_2^2\).
\end{proof}

\subsection{Certificate via Weyl's inequality}
Let \(A\) and \(\Delta\) be Hermitian. By Weyl's inequality, for each \(i\),
\[
  \abs{\lambda_i(A+\Delta)-\lambda_i(A)} \le \norm{\Delta}_2.
\]
Therefore the change in any gap satisfies
\begin{equation}
\label{eq:gap-bound}
  \abs{g(A+\Delta) - g(A)} \le 2\,\norm{\Delta}_2.
\end{equation}
Thus, constraining \(\norm{\Delta}_2\le \varepsilon\) certifies that the gap changes by at most \(2\varepsilon\).

\begin{proposition}[Gap-control certificate]
\label{prop:gap-certificate}
If the SCN produces a perturbation \(\Delta\) with \(\norm{\Delta}_2\le \varepsilon\), then the spectral gap deviation is bounded by \eqref{eq:gap-bound}. 
\end{proposition}

\subsection{SCN architecture and objective}
We define a feature map \(\phi\) collecting spectral statistics:
\[
  \phi(A,\hat g) = \big[g(A),\, \hat g,\, \operatorname{mean}(\lambda),\, \operatorname{std}(\lambda),\, \min \lambda,\, \max \lambda,\, \operatorname{mean}(\Delta\lambda),\, \operatorname{std}(\Delta\lambda)\big] \in \R^{s}.
\]
A small network \(F_\theta:\R^{s}\to\R^{m}\) outputs a control vector \(w\), and a differentiable map \(w\mapsto \Delta(w)\) yields the proposed perturbation. 
Training minimizes
\[
  \mathcal{L}(\theta) 
  = \underbrACE/SCN{\big(g(A+\Delta(w_\theta))-\hat g\big)^2}_{\text{gap tracking}}
  + \lambda \underbrACE/SCN{\max\{0,\norm{\Delta(w_\theta)}_2-\varepsilon\}^2}_{\text{certificate penalty}},
\]
optionally with regularizers on \(w\) or structure in \(\Delta\).

\paragraph{Minimal code excerpt.}
\begin{lstlisting}[language=Python,caption={Spectral Control Network (unitary block).}]
import torch, torch.nn as nn
import numpy as np

class UnitaryBlock(nn.Module):
    def __init__(self, dim: int):
        super().__init__()
        self.A = nn.Parameter(torch.randn(dim, dim) * 0.01)
    def forward(self, x: torch.Tensor) -> torch.Tensor:
        S = self.A - self.A.T
        U = torch.matrix_exp(S)    # orthogonal in R, unitary in C
        return x @ U

class SpectralControlNetwork(nn.Module):
    def __init__(self, state_dim=8, hidden_dim=32, control_dim=8):
        super().__init__()
        self.inp = nn.Linear(state_dim, hidden_dim)
        self.unitary = UnitaryBlock(hidden_dim)
        self.out = nn.Linear(hidden_dim, control_dim)
    def forward(self, state: torch.Tensor) -> torch.Tensor:
        h = torch.tanh(self.inp(state))
        h = self.unitary(h)
        return torch.tanh(self.out(h))
\end{lstlisting}

% =========================================================
\section{Validation and Scientific Tests}
\label{sec:validation}

\subsection{Arithmetic tests}
\begin{itemize}[leftmargin=*,itemsep=2pt]
  \item \textbf{Hecke multiplicativity:} verify \(\norm{T_mT_n - T_{mn}}\le \varepsilon\) for coprime \((m,n)\).
  \item \textbf{Ramanujan--Petersson:} verify \(\abs{a_p}\le 2p^{(k-1)/2}+\varepsilon\) for a set of primes \(p\nmid N\).
\end{itemize}

\subsection{Statistical distribution checks}
When comparing empirical spectral statistics to a reference distribution, we employ the Kolmogorov--Smirnov two-sample test and report the statistic and \(p\)-value; we \emph{do not} claim agreement without explicit thresholds and power discussion.

\subsection{Reproducibility}
To ensure determinism:
\begin{enumerate}[leftmargin=*,itemsep=2pt]
  \item Fix random seeds across libraries.
  \item Pin package versions.
  \item Provide exact scripts for each table/figure.
\end{enumerate}

\subsection{Illustrative test excerpts}
\begin{lstlisting}[language=Python,caption={Ramanujan--Petersson bound test.}]
from ACE/SCN.core.hecke_core import hecke_eigenvalues
def verify_rp(evals, p, k, tol=1e-8):
    bound = 2 * (p ** ((k-1)/2))
    return all(abs(complex(ev)) <= bound + tol for ev in evals)

def test_rp_level1_weight12():
    primes = [2,3,5,7,11,13,17,19,23]
    for p in primes:
        evals = hecke_eigenvalues(level=1, weight=12, n=p)
        assert verify_rp(evals, p, 12)
\end{lstlisting}

% =========================================================
\section{Repository Layout and Workflow}
\label{sec:repo}

\subsection{Minimal structure}
\begin{lstlisting}[language=bash,caption={Directory layout.}]
ACE/SCN-Spectral-Control/
  ACE/SCN/
    core/           # Hecke operators (exact)
    control/        # Spectral Control Network
    validation/     # Bound checks, stats, units
  tests/
    scientific/     # Theorem-backed tests
    unit/           # Lightweight unit tests
  examples/         # Minimal end-to-end scripts
  requirements.txt
  README.md
\end{lstlisting}

\subsection{Quickstart}
\begin{enumerate}[leftmargin=*,itemsep=2pt]
  \item Install a computer algebra system that provides \(T_n\) on \(S_k(\Gamma_0(N))\).
  \item \texttt{pip install -r requirements.txt}.
  \item Run \texttt{python tests/scientific/test\_ramanujan\_bounds.py}.
  \item Explore \texttt{examples/basic\_hecke.py} and \texttt{examples/scn\_demo.py}.
\end{enumerate}

% =========================================================
\section{Limitations and Scope}
\label{sec:limitations}
ACE/SCN currently focuses on:
\begin{itemize}[leftmargin=*,itemsep=2pt]
  \item Classical cusp forms \(S_k(\Gamma_0(N))\) and prime-power Hecke operators.
  \item Hermitian/symmetric operators for gap control and certification via \eqref{eq:gap-bound}.
\end{itemize}
We do not claim:
\begin{itemize}[leftmargin=*,itemsep=2pt]
  \item New theorems in number theory.
  \item Quantum computation or hardware-level models.
  \item Production infrastructure beyond research reproducibility.
\end{itemize}

% =========================================================
\section{Roadmap}
\label{sec:roadmap}
\begin{enumerate}[leftmargin=*,itemsep=2pt]
  \item Extend arithmetic functionality to new levels \(N\), weights \(k\), and Atkin--Lehner operators; add old/newform decompositions.
  \item Enrich SCN with structured perturbations \(\Delta(w)\) aligned with known operator symmetries.
  \item Add property-based tests for spectral gap shaping; benchmark against curated modular-form datasets.
  \item Package reproducible containers for exact-environment runs.
\end{enumerate}

% =========================================================
\section*{Acknowledgments}
We thank the community for careful critiques that motivated this revision: replacing superficial terminology with explicit constructions and falsifiable tests.

% =========================================================
\appendix

\section{Educational Fallback (Optional)}
\label{app:fallback}
For environments without computer algebra, an \emph{educational stub} may synthesize small symmetric matrices to exercise the SCN and the validation harness.
These stubs are clearly marked and never used to make arithmetic claims.

\section{Environment and Seeding}
\label{app:env}
We fix seeds in NumPy/PyTorch; versions are pinned in \texttt{requirements.txt}. For exact arithmetic, we rely on the CAS implementation of Hecke operators.

% =========================================================
\begin{thebibliography}{9}

\bibitem{diamondshurman}
F.~Diamond and J.~Shurman,
\newblock \emph{A First Course in Modular Forms},
\newblock Springer, 2005.

\bibitem{serre}
J.-P.~Serre,
\newblock \emph{A Course in Arithmetic},
\newblock Springer, 1973.

\bibitem{iwaniecluo}
H.~Iwaniec,
\newblock \emph{Topics in Classical Automorphic Forms},
\newblock American Mathematical Society, 1997.

\end{thebibliography}

\end{document}



